% Einbindung nach dem vollständig bewiesenen Strukturrekursionssatz.
\subsection{Der Baumrekursor: Blätter ersetzen, Ergebnisse verknüpfen}

Der Baumrekursor führt zwei festgelegte Schritte aus: Er ersetzt ein Blatt
mit Beschriftung \(a\) durch \(f(a)\), und er verknüpft die Ergebnisse der
beiden Teilbäume mit \(g\). Der Strukturrekursionssatz liefert genau eine
Funktion mit diesen Eigenschaften.

\Needspace{16\baselineskip}
\FormulaDefDeltaK[Baumrekursor]{%
  \begin{aligned}[t]
    &f\colon A\to X\dsep g\colon X\times X\to X\vdash{}\\[-2pt]
    &\TreeRecursor{A}{X}{f}{g}\coloneqq
      \iota E\colon\BracketTreeSet A\to X\,\Bigl(\\[-2pt]
    &\quad\forall a\in A\,E(\LeafTree A a)=f(a)\land{}\\[-2pt]
    &\quad\forall S,T\in\BracketTreeSet A\,
      E(\NodeTree A S T)=g(E(S),E(T))\Bigr)
  \end{aligned}%
}{TreeRecursorDef}{%
  \DeltaRow{Mengen}{A\dsep X\dsep a\dsep S\dsep T}
  \DeltaRow{Funktionen}{f\dsep g\dsep E}
  \DeltaRow{\textbf{Neues Symbol}}{}
  \DeltaPrem{Baumrekursor}{\TreeRecursor{A}{X}{f}{g}}
}

Für die folgenden Aussagen und Beweise wird
\(R\coloneqq\TreeRecursor{A}{X}{f}{g}\) abgekürzt.

\FormulaThmDeltaK[Charakterisierung des Baumrekursors]{%
  \begin{aligned}[t]
    &f\colon A\to X\dsep g\colon X\times X\to X\vdash{}\\[-2pt]
    &R\colon\BracketTreeSet A\to X\land\Bigl(\\[-2pt]
    &\quad\forall a\in A\,R(\LeafTree A a)=f(a)\land{}\\[-2pt]
    &\quad\forall S,T\in\BracketTreeSet A\,
      R(\NodeTree A S T)=g(R(S),R(T))\Bigr)
  \end{aligned}%
}{TreeRecursorSpecification}{%
  \DeltaRow{Mengen}{A\dsep X\dsep a\dsep S\dsep T}
  \DeltaRow{Funktionen}{f\dsep g\dsep E\dsep R}
  \DeltaRow{Abkürzung}{R\coloneqq\TreeRecursor{A}{X}{f}{g}}
}
\begin{tabproofwide}
  \proofstepwidestar[1]{f\colon A\to X}{\rA}
  \proofstepwidestar[2]{g\colon X\times X\to X}{\rA}
  \proofstepwidestar[1,2]{%
    \begin{aligned}[t]
      &\exists!E\colon\BracketTreeSet A\to X\,\Bigl(\\[-2pt]
      &\quad\forall a\in A\,E(\LeafTree A a)=f(a)\land{}\\[-2pt]
      &\quad\forall S,T\in\BracketTreeSet A\,
        E(\NodeTree A S T)\\[-2pt]
      &\qquad=g(E(S),E(T))\Bigr)
    \end{aligned}}{%
    \FormulaRefAuto{FullPlanarBinaryTreeRecursion}[thm]{1,2}}
  \proofstepwidestar[1,2]{%
    \begin{aligned}[t]
      &R\colon\BracketTreeSet A\to X\land\Bigl(\\[-2pt]
      &\quad\forall a\in A\,R(\LeafTree A a)=f(a)\land{}\\[-2pt]
      &\quad\forall S,T\in\BracketTreeSet A\,
        R(\NodeTree A S T)\\[-2pt]
      &\qquad=g(R(S),R(T))\Bigr)
    \end{aligned}}{%
    \FormulaRefAuto{TreeRecursorDef}[def]{1,2,3}}
\end{tabproofwide}

\Needspace{20\baselineskip}
\FormulaThmDeltaK[Typisierung, Blattregel und Knotenregel]{%
  \begin{aligned}[t]
    &\text{(i)}\;f\colon A\to X\dsep g\colon X\times X\to X\vdash
      R\colon\BracketTreeSet A\to X,\\[-2pt]
    &\text{(ii)}\;f\colon A\to X\dsep g\colon X\times X\to X\dsep
      a\in A\vdash R(\LeafTree A a)=f(a),\\[-2pt]
    &\text{(iii)}\;f\colon A\to X\dsep g\colon X\times X\to X\dsep
      S,T\in\BracketTreeSet A\vdash{}\\[-2pt]
    &\qquad R(\NodeTree A S T)=g(R(S),R(T)).
  \end{aligned}%
}{TreeRecursorEquations}{%
  \DeltaRow{Mengen}{A\dsep X\dsep a\dsep S\dsep T}
  \DeltaRow{Funktionen}{f\dsep g\dsep R}
  \DeltaRow{Abkürzung}{R\coloneqq\TreeRecursor{A}{X}{f}{g}}
}
\begin{tabproofsplitwide}
  \Needspace{10\baselineskip}
  \proofpartwideindR[Typisierung]{%
    f\colon A\to X\dsep g\colon X\times X\to X\vdash
      R\colon\BracketTreeSet A\to X
  }{f\colon A\to X\dsep g\colon X\times X\to X\vdash
      R\colon\BracketTreeSet A\to X}[TreeRecursorTyping]
    \proofstepwidestar[1]{f\colon A\to X}{\rA}
    \proofstepwidestar[2]{g\colon X\times X\to X}{\rA}
    \proofstepwidestar[1,2]{%
      \begin{aligned}[t]
        &R\colon\BracketTreeSet A\to X\land\Bigl(\\[-2pt]
        &\quad\forall a\in A\,R(\LeafTree A a)=f(a)\land{}\\[-2pt]
        &\quad\forall S,T\in\BracketTreeSet A\,R(\NodeTree A S T)\\[-2pt]
        &\qquad=g(R(S),R(T))\Bigr)
      \end{aligned}}{\FormulaRefAuto{TreeRecursorSpecification}[thm]{1,2}}
    \proofstepwidestar[1,2]{%
      R\colon\BracketTreeSet A\to X}{%
      \rAEa{3}}
  \closeproofpartwideindR
  \Needspace{10\baselineskip}
  \proofpartwideindR[Blattregel]{%
    f\colon A\to X\dsep g\colon X\times X\to X\dsep a\in A\vdash
      R(\LeafTree A a)=f(a)
  }{f\colon A\to X\dsep g\colon X\times X\to X\dsep a\in A\vdash
      R(\LeafTree A a)=f(a)}[TreeRecursorLeafEquation]
    \proofstepwidestar[1]{f\colon A\to X}{\rA}
    \proofstepwidestar[2]{g\colon X\times X\to X}{\rA}
    \proofstepwidestar[3]{a\in A}{\rA}
    \proofstepwidestar[1,2]{%
      \begin{aligned}[t]
        &R\colon\BracketTreeSet A\to X\land\Bigl(\\[-2pt]
        &\quad\forall a\in A\,R(\LeafTree A a)=f(a)\land{}\\[-2pt]
        &\quad\forall S,T\in\BracketTreeSet A\,R(\NodeTree A S T)\\[-2pt]
        &\qquad=g(R(S),R(T))\Bigr)
      \end{aligned}}{\FormulaRefAuto{TreeRecursorSpecification}[thm]{1,2}}
    \proofstepwidestar[1,2]{%
      \forall a\in A\,R(\LeafTree A a)=f(a)}{%
      \rAEa{\rAEb{4}}}
    \proofstepwidestar[1,2]{a\in A\rightarrow
      R(\LeafTree A a)=f(a)}{\rUE{5}}
    \proofstepwidestar[1,2,3]{R(\LeafTree A a)=f(a)}{\rRE{3,6}}
  \closeproofpartwideindR
  \Needspace{10\baselineskip}
  \proofpartwideindR[Knotenregel]{%
    \begin{aligned}[t]
      &f\colon A\to X\dsep g\colon X\times X\to X\dsep
        S,T\in\BracketTreeSet A\vdash{}\\[-2pt]
      &R(\NodeTree A S T)=g(R(S),R(T))
    \end{aligned}
  }{\begin{aligned}[t]
      &f\colon A\to X\dsep g\colon X\times X\to X\dsep
        S,T\in\BracketTreeSet A\vdash{}\\[-2pt]
      &R(\NodeTree A S T)=g(R(S),R(T))
    \end{aligned}}[TreeRecursorNodeEquation]
    \proofstepwidestar[1]{f\colon A\to X}{\rA}
    \proofstepwidestar[2]{g\colon X\times X\to X}{\rA}
    \proofstepwidestar[3]{S\in\BracketTreeSet A}{\rA}
    \proofstepwidestar[4]{T\in\BracketTreeSet A}{\rA}
    \proofstepwidestar[1,2]{%
      \begin{aligned}[t]
        &R\colon\BracketTreeSet A\to X\land\Bigl(\\[-2pt]
        &\quad\forall a\in A\,R(\LeafTree A a)=f(a)\land{}\\[-2pt]
        &\quad\forall S,T\in\BracketTreeSet A\,R(\NodeTree A S T)\\[-2pt]
        &\qquad=g(R(S),R(T))\Bigr)
      \end{aligned}}{\FormulaRefAuto{TreeRecursorSpecification}[thm]{1,2}}
    \proofstepwidestar[1,2]{%
      \forall S,T\in\BracketTreeSet A\,
        R(\NodeTree A S T)=g(R(S),R(T))}{%
      \rAEb{\rAEb{5}}}
    \proofstepwidestar[1,2]{%
      S\in\BracketTreeSet A\rightarrow
      \forall T\in\BracketTreeSet A\,
        R(\NodeTree A S T)=g(R(S),R(T))}{\rUE{6}}
    \proofstepwidestar[1,2,3]{%
      \forall T\in\BracketTreeSet A\,
        R(\NodeTree A S T)=g(R(S),R(T))}{\rRE{3,7}}
    \proofstepwidestar[1,2,3]{%
      T\in\BracketTreeSet A\rightarrow
        R(\NodeTree A S T)=g(R(S),R(T))}{\rUE{8}}
    \proofstepwidestar[1,2,3,4]{%
      R(\NodeTree A S T)=g(R(S),R(T))}{\rRE{4,9}}
  \closeproofpartwideindR
\end{tabproofsplitwide}

Die nächste Regel identifiziert eine bereits gegebene Funktion: Es genügt,
ihre Typisierung sowie ihre Blatt- und Knotenregel nachzuweisen.

\FormulaThmDeltaK[Identifikationsregel für den Baumrekursor]{%
  \begin{aligned}[t]
    &f\colon A\to X\dsep g\colon X\times X\to X\dsep
      E\colon\BracketTreeSet A\to X\dsep{}\\[-2pt]
    &\forall a\in A\,E(\LeafTree A a)=f(a)\dsep{}\\[-2pt]
    &\forall S,T\in\BracketTreeSet A\,
      E(\NodeTree A S T)=g(E(S),E(T))\vdash E=R
  \end{aligned}%
}{TreeRecursorIdentification}{%
  \DeltaRow{Mengen}{A\dsep X\dsep a\dsep S\dsep T}
  \DeltaRow{Funktionen}{f\dsep g\dsep E\dsep R}
  \DeltaRow{Abkürzung}{R\coloneqq\TreeRecursor{A}{X}{f}{g}}
}
\begin{tabproofwide}
  \proofstepwidestar[1]{f\colon A\to X}{\rA}
  \proofstepwidestar[2]{g\colon X\times X\to X}{\rA}
  \proofstepwidestar[3]{E\colon\BracketTreeSet A\to X}{\rA}
  \proofstepwidestar[4]{\forall a\in A\,E(\LeafTree A a)=f(a)}{\rA}
  \proofstepwidestar[5]{%
    \forall S,T\in\BracketTreeSet A\,
      E(\NodeTree A S T)=g(E(S),E(T))}{\rA}
  \proofstepwidestar[1,2]{%
    \begin{aligned}[t]
      &R\colon\BracketTreeSet A\to X\land\Bigl(\\[-2pt]
      &\quad\forall a\in A\,R(\LeafTree A a)=f(a)\land{}\\[-2pt]
      &\quad\forall S,T\in\BracketTreeSet A\,
        R(\NodeTree A S T)\\[-2pt]
      &\qquad=g(R(S),R(T))\Bigr)
    \end{aligned}}{\FormulaRefAuto{TreeRecursorSpecification}[thm]{1,2}}
  \proofstepwidestar[1,2]{R\colon\BracketTreeSet A\to X}{\rAEa{6}}
  \proofstepwidestar[1,2]{%
    \forall a\in A\,R(\LeafTree A a)=f(a)}{\rAEa{\rAEb{6}}}
  \proofstepwidestar[1,2]{%
    \forall S,T\in\BracketTreeSet A\,
      R(\NodeTree A S T)=g(R(S),R(T))}{\rAEb{\rAEb{6}}}
  \proofstepwidestar[1,2,3]{E,R\colon\BracketTreeSet A\to X}{\rAI{3,7}}
  \proofstepwidestar[1,2,3,4,5]{E=R}{%
    \FormulaRefAuto{\exists! x\,P(x),\; P(a),\; P(b) \vdash a = b}[thm]{%
      \FormulaRefAuto{FullPlanarBinaryTreeRecursion}[thm]{1,2},
      \rAI{3,\rAI{4,5}},6}}
\end{tabproofwide}

\subsection{Rekursion mit einer binären Operation}

Bei einer binären Operation \(\star\) wird ihre bereits definierte
zugehörige Funktion eingesetzt. Wir schreiben hier
\(R_\star\coloneqq
\TreeRecursor{A}{X}{f}{\BinaryOpFunction{X}{\star}}\).
Die Knotenregel verwendet dann unmittelbar das Operationszeichen.

\FormulaThmDeltaK[Baumrekursor mit binärer Operation]{%
  \begin{aligned}[t]
    &\text{(i)}\;f\colon A\to X\dsep\mathsf{BinOp}(X,\star)\vdash
      R_\star\colon\BracketTreeSet A\to X,\\[-2pt]
    &\text{(ii)}\;f\colon A\to X\dsep\mathsf{BinOp}(X,\star)\dsep
      a\in A\vdash R_\star(\LeafTree A a)=f(a),\\[-2pt]
    &\text{(iii)}\;f\colon A\to X\dsep\mathsf{BinOp}(X,\star)\dsep
      S,T\in\BracketTreeSet A\vdash{}\\[-2pt]
    &\qquad R_\star(\NodeTree A S T)=R_\star(S)\star R_\star(T).
  \end{aligned}%
}{TreeRecursorBinaryOperationEquations}{%
  \DeltaRow{Mengen}{A\dsep X\dsep a\dsep S\dsep T}
  \DeltaRow{Funktionen}{f\dsep\star\dsep R_\star}
  \DeltaRow{Abkürzung}{R_\star\coloneqq
    \TreeRecursor{A}{X}{f}{\BinaryOpFunction{X}{\star}}}
}
\begin{tabproofsplitwide}
  \Needspace{10\baselineskip}
  \proofpartwideindR[Typisierung]{%
    f\colon A\to X\dsep\mathsf{BinOp}(X,\star)\vdash
      R_\star\colon\BracketTreeSet A\to X
  }{f\colon A\to X\dsep\mathsf{BinOp}(X,\star)\vdash
      R_\star\colon\BracketTreeSet A\to X}[TreeRecursorBinaryOperationTyping]
    \proofstepwidestar[1]{f\colon A\to X}{\rA}
    \proofstepwidestar[2]{\mathsf{BinOp}(X,\star)}{\rA}
    \proofstepwidestar[2]{%
      \BinaryOpFunction{X}{\star}\colon X\times X\to X}{%
      \FormulaRefAuto{BinaryOperationFunctionDef}[def]{2}}
    \proofstepwidestar[1,2]{R_\star\colon\BracketTreeSet A\to X}{%
      \FormulaRefAuto{TreeRecursorTyping}[thm]{1,3}}
  \closeproofpartwideindR
  \Needspace{10\baselineskip}
  \proofpartwideindR[Blattregel]{%
    f\colon A\to X\dsep\mathsf{BinOp}(X,\star)\dsep a\in A\vdash
      R_\star(\LeafTree A a)=f(a)
  }{f\colon A\to X\dsep\mathsf{BinOp}(X,\star)\dsep a\in A\vdash
      R_\star(\LeafTree A a)=f(a)}[TreeRecursorBinaryOperationLeafEquation]
    \proofstepwidestar[1]{f\colon A\to X}{\rA}
    \proofstepwidestar[2]{\mathsf{BinOp}(X,\star)}{\rA}
    \proofstepwidestar[3]{a\in A}{\rA}
    \proofstepwidestar[2]{%
      \BinaryOpFunction{X}{\star}\colon X\times X\to X}{%
      \FormulaRefAuto{BinaryOperationFunctionDef}[def]{2}}
    \proofstepwidestar[1,2,3]{R_\star(\LeafTree A a)=f(a)}{%
      \FormulaRefAuto{TreeRecursorLeafEquation}[thm]{1,4,3}}
  \closeproofpartwideindR
  \Needspace{10\baselineskip}
  \proofpartwideindR[Knotenregel]{%
    \begin{aligned}[t]
      &f\colon A\to X\dsep\mathsf{BinOp}(X,\star)\dsep
        S,T\in\BracketTreeSet A\vdash{}\\[-2pt]
      &R_\star(\NodeTree A S T)=R_\star(S)\star R_\star(T)
    \end{aligned}
  }{\begin{aligned}[t]
      &f\colon A\to X\dsep\mathsf{BinOp}(X,\star)\dsep
        S,T\in\BracketTreeSet A\vdash{}\\[-2pt]
      &R_\star(\NodeTree A S T)=R_\star(S)\star R_\star(T)
    \end{aligned}}[TreeRecursorBinaryOperationNodeEquation]
    \proofstepwidestar[1]{f\colon A\to X}{\rA}
    \proofstepwidestar[2]{\mathsf{BinOp}(X,\star)}{\rA}
    \proofstepwidestar[3]{S\in\BracketTreeSet A}{\rA}
    \proofstepwidestar[4]{T\in\BracketTreeSet A}{\rA}
    \proofstepwidestar[2]{%
      \BinaryOpFunction{X}{\star}\colon X\times X\to X}{%
      \FormulaRefAuto{BinaryOperationFunctionDef}[def]{2}}
    \proofstepwidestar[1,2]{R_\star\colon\BracketTreeSet A\to X}{%
      \FormulaRefAuto{TreeRecursorTyping}[thm]{1,5}}
    \proofstepwidestar[1,2,3]{R_\star(S)\in X}{%
      \FormulaRefAuto{F\colon A\to B\dsep x\in A\vdash F(x)\in B}[thm]{6,3}}
    \proofstepwidestar[1,2,4]{R_\star(T)\in X}{%
      \FormulaRefAuto{F\colon A\to B\dsep x\in A\vdash F(x)\in B}[thm]{6,4}}
    \proofstepwidestar[1,2,3,4]{%
      R_\star(\NodeTree A S T)=
        \BinaryOpFunction{X}{\star}(R_\star(S),R_\star(T))}{%
      \FormulaRefAuto{TreeRecursorNodeEquation}[thm]{1,5,\rAI{3,4}}}
    \proofstepwidestar[1,2,3,4]{%
      \BinaryOpFunction{X}{\star}(R_\star(S),R_\star(T))=
        R_\star(S)\star R_\star(T)}{%
      \FormulaRefAuto{BinaryOperationFunctionDef}[def]{2,7,8}}
    \proofstepwidestar[1,2,3,4]{%
      R_\star(\NodeTree A S T)=R_\star(S)\star R_\star(T)}{%
      \FormulaRefAuto{a=b\dsep b=c\vdash a=c}[thm]{9,10}}
  \closeproofpartwideindR
\end{tabproofsplitwide}

\FormulaThmDeltaK[Identifikation bei einer binären Operation]{%
  \begin{aligned}[t]
    &f\colon A\to X\dsep\mathsf{BinOp}(X,\star)\dsep
      E\colon\BracketTreeSet A\to X\dsep{}\\[-2pt]
    &\forall a\in A\,E(\LeafTree A a)=f(a)\dsep{}\\[-2pt]
    &\forall S,T\in\BracketTreeSet A\,
      E(\NodeTree A S T)=E(S)\star E(T)\vdash E=R_\star
  \end{aligned}%
}{TreeRecursorBinaryOperationIdentification}{%
  \DeltaRow{Mengen}{A\dsep X\dsep a\dsep S\dsep T}
  \DeltaRow{Funktionen}{f\dsep\star\dsep E\dsep R_\star}
  \DeltaRow{Abkürzung}{R_\star\coloneqq
    \TreeRecursor{A}{X}{f}{\BinaryOpFunction{X}{\star}}}
}
\begin{tabproofwide}
  \proofstepwidestar[1]{f\colon A\to X}{\rA}
  \proofstepwidestar[2]{\mathsf{BinOp}(X,\star)}{\rA}
  \proofstepwidestar[3]{E\colon\BracketTreeSet A\to X}{\rA}
  \proofstepwidestar[4]{\forall a\in A\,E(\LeafTree A a)=f(a)}{\rA}
  \proofstepwidestar[5]{%
    \forall S,T\in\BracketTreeSet A\,
      E(\NodeTree A S T)=E(S)\star E(T)}{\rA}
  \proofstepwidestar[2]{%
    \BinaryOpFunction{X}{\star}\colon X\times X\to X}{%
    \FormulaRefAuto{BinaryOperationFunctionDef}[def]{2}}
  \proofstepwidestar[7]{S\in\BracketTreeSet A}{\rA}
  \proofstepwidestar[8]{T\in\BracketTreeSet A}{\rA}
  \proofstepwidestar[3,7]{E(S)\in X}{%
    \FormulaRefAuto{F\colon A\to B\dsep x\in A\vdash F(x)\in B}[thm]{3,7}}
  \proofstepwidestar[3,8]{E(T)\in X}{%
    \FormulaRefAuto{F\colon A\to B\dsep x\in A\vdash F(x)\in B}[thm]{3,8}}
  \proofstepwidestar[5,7,8]{%
    E(\NodeTree A S T)=E(S)\star E(T)}{%
    \rRE{8,\rUE{\rRE{7,\rUE{5}}}}}
  \proofstepwidestar[2,3,7,8]{%
    \BinaryOpFunction{X}{\star}(E(S),E(T))=E(S)\star E(T)}{%
    \FormulaRefAuto{BinaryOperationFunctionDef}[def]{2,9,10}}
  \proofstepwidestar[2,3,5,7,8]{%
    E(\NodeTree A S T)=\BinaryOpFunction{X}{\star}(E(S),E(T))}{%
    \FormulaRefAuto{a=b\dsep c=b\vdash a=c}[thm]{11,12}}
  \proofstepwidestar[2,3,5]{%
    \forall S,T\in\BracketTreeSet A\,
      E(\NodeTree A S T)=\BinaryOpFunction{X}{\star}(E(S),E(T))}{%
    \rUI{\rRI{7,\rUI{\rRI{8,13}}}}}
  \proofstepwidestar[1,2,3,4,5]{E=R_\star}{%
    \FormulaRefAuto{TreeRecursorIdentification}[thm]{1,6,3,4,14}}
\end{tabproofwide}
